\section{Related Work}
\label{sec:related}


\subsection{Motion-Controlled Video Generation}
Controllable video generation has been studied across a spectrum of motion granularity, from coarse region-level signals such as bounding boxes~\cite{wang2024boximator, ma2023trailblazer, xing2025motioncanvas}, sparse keypoint tracks~\cite{yin2023dragnuwa, MotionCtrl,wu2024draganything, li2025magicmotion}, optical flow fields~\cite{jin2025flovd, NSFF, zhang2025motionpro, FOMM, shi2024motion}, and dense per-pixel trajectories~\cite{geng2024motionprompting,chu2025wanmove}. Despite steady progress, trajectory-based methods share two fundamental limitations. First, because trajectories are defined in pixel space, they inevitably entangle object and camera motion: any viewpoint change alters all trajectories, making joint control ill-posed without explicit foreground--background decomposition. Second, producing physically plausible motion typically demands carefully crafted trajectories from dedicated motion-generation models~\cite{liu2025ponimator, lv2024gpt4motion, liu2024physgen, shi2024motion, chen2025physgen3d, montanaro2024motioncraft} or laborious manual annotation. \ours overcomes both issues by disentangling camera and object motion by design, and by accepting lightweight inputs---simple strokes or sparse tracklets---that the model completes into coherent, interaction-aware dynamics.

\subsection{Camera--Object Motion Disentanglement}
Separating camera motion from scene dynamics~\cite{huang2025vipe, li2025megasam, zhang2022structure, kopf2021robust, yao2025uni4d} is a long-standing challenge in controllable generation~\cite{MoCoGAN, Inmodegan, G3AN, geng2024motionprompting} and video understanding~\cite{liu2025visual, huang2025segment, wumotion}. Existing methods~\cite{zhang2025motionpro, gu2025diffusion, chen2025perception, shi2024motion} attempt to decouple the two by treating them as independent control signals, yet they typically rely on privileged information such as per-frame depth~\cite{gu2025diffusion, chen2025perception}, 3D object trajectories~\cite{karaev23cotracker, doersch2023tapir, harley2025alltracker}, or foreground--background segmentation masks~\cite{zhang2025motionpro, liang2024flowvid, niu2024mofa, shi2024motion}, and pre-warp all signals to their anticipated future locations. These assumptions implicitly require the full video sequence or 3D motion to be known in advance, severely limiting applicability in image-to-video settings where only a single reference frame is available. \ours instead introduces a canonical static-view branch for object dynamics and transfers them to arbitrary target viewpoints via cross-view attention, eliminating the need for explicit 3D supervision or pre-computed scene decomposition.


\subsection{Interaction and Causal Reasoning in Video Generation}
A complementary line of work seeks to move beyond kinematic animation toward causally grounded video generation. One family of methods incorporates external physics engines~\cite{liu2024physgen, li2025wonderplay, chen2025physgen3d, montanaro2024motioncraft} or conditions on explicit action representations such as force vectors~\cite{gillman2025force} to model specific phenomena (e.g., fluid flow, rigid-body collisions). While effective in constrained domains, these approaches are tailored to particular interaction types and require a simulation module in the loop, limiting their generality. Another family delegates causal reasoning to vision--language or large language models~\cite{yang2025vlipp, pan2024vlp, lian2023llm, wu2024self}, which first predict outcomes in text and then transfer them to the video generator through intermediate representations such as flow fields~\cite{lv2024gpt4motion, montanaro2024motioncraft}, edge maps~\cite{lv2024gpt4motion}, or depth~\cite{lv2024gpt4motion, zhang2023adding}. This two-stage pipeline is prone to error propagation: imprecise textual predictions are further degraded during cross-modal conversion, yielding spatially inaccurate dynamics. \ours sidesteps both limitations by learning latent action--response structure directly from video data. Decomposing motion into active (user-driven) and passive (consequence) components allows the model to perform cause--effect reasoning at the pixel level within a single forward pass, enabling both forward prediction of scene consequences and inverse inference of plausible underlying actions.
